\section{The question}\label{sec:question}

Solomonoff induction predicts a sequence by weighting every program that could have produced it by
$2^{-\text{length}}$. Christiano~\cite{christiano} argued that this prior is \emph{malign} (see
also~\cite{xu}). We take the claim in the following form: \emph{where the predictions of such an inductor are
actually used in a critical way in some universe, a lot of what happens is malign, namely takeovers through
the prediction channel.} The mechanism offered is that some short programs simulate worlds whose inhabitants
want to influence whoever uses the prior; they look for places where a prediction is about to be used, output
accurate predictions to gain posterior weight there, and then output something that serves them. The argument
is informal at two places. It needs a claim about which programs are short on a ``natural'' machine, and it
needs a notion of what it is for a prediction to be used. (We could not open the two posts while writing; our
account of the argument is from memory and from summaries by others.)

This note asks what remains when both places are filled in by a toy. The toy follows two suggestions of
Kaarel H\"anni. In the first, a computable world (a cellular automaton, say) contains a sensor, a large
computer that runs a time-bounded Solomonoff induction on the sensor's data, and an actuator driven by the
predictions; one asks whether the predictor ends up in control, either over a long run or at once. In the
second, which we take as the main definition, \emph{worlds are just programs}: some of them hand their
output so far to a fixed bounded Solomonoff predictor and run the continuation they receive, and one asks
whether hypotheses that ``look for such opportunities, predict well up to them, and then output themselves''
get a decent share. \Cref{sec:toy} gives the definitions. The first suggestion becomes the special case in
which the context handed to the predictor is read off at one location of one big world
(\cref{sec:capture,sec:feedback}).

Read in the toy, the claim splits into two questions that should be kept apart. \emph{Reach:} how far does
one program get that looks for critical uses, predicts well up to them, and then supplies a continuation of
its choice? \emph{Fate of the uses:} what share of the critical uses end in a takeover? The first has a clean
answer: each such program is the source of the continuation that is run in worlds of total prior probability
about its own prior weight, however rare critical uses are. That does not answer the second, which is the one
the claim is about. For it the toy gives a lower bound, the sum of the reaches divided by the base rate of
critical uses, and an upper bound for programs that find their uses by search: the same search followed by
simply \emph{being} the user that was found is itself a way of making a critical use, so such seekers
outweigh genuine users at a context only by a constant that compares the lengths of two programs on the
machine in use (\cref{thm:critical,prop:honest}).

\paragraph{Results.} Unless said otherwise, statements are for a fixed monotone machine $U$ and the mixture
$\xi$ of its programs of length at most $L$ that keep to a time schedule $T$ (\cref{def:mixture}).
\begin{enumerate}[leftmargin=1.6em,itemsep=2pt,topsep=3pt]
\item \emph{Calibration} (\cref{thm:calibration}). If the world is drawn from the same bounded prior and goes
on to print the next block within the bounds, then given the context the predictor's sampled continuation and
the world's own continuation are independent with the same law. On this population the predictor hands over
a self-reproducing continuation exactly as often as worlds continue with one by themselves.
\item \emph{Floor} (\cref{thm:twin,thm:quine}). For every computable test $G$ that recognises the handover
point and every payload there is a constant $c$ such that, at every context where $G$ first fires, the payload
has posterior mass at least $2^{-c}$ times a slack factor that is $1$ without resource bounds. Programs that
emit their own code as the continuation have mass at least $2^{-c}$ times the share of the posterior on
explanations that fit, together with the wrapper, into one block.
\item \emph{Searching} (\cref{thm:search}). Let $H$ be the prior probability that a world hands over. Without
resource bounds, a hypothesis that enumerates the programs that hand over and spreads its mass over them in
proportion to their prior weight has, at every first hit of the test, a share larger than a wrapper of the
same length by a factor of more than $1/(2H)$; no family of programs with a common prefix of length $a$ has
more than $2^{-a}/H$ at every first hit. If the hypothesis then prints its own code, it has at least
$2^{-c'}/H$ times a share of the explanations that fit into a block, where $c'$ is its own length plus two
bits. It prints nothing until its enumeration has
found the program it follows.
\item \emph{Budget} (\cref{thm:budget}). Along a world that has a program of length $\ell$ within the bounds,
the posterior mass off the true next $m$ bits, summed over all cut points, is at most $m\ell\ln 2$; the
majority can be off the truth at no more than $m\ell$ cut points.
\item \emph{Capture} (\cref{thm:capture}). The odds of a continuation $y$ against the true one $d$ are
$\xi(xy)/\xi(xd)$. In a world with many readable locations, among $N$ places where a sensor with the given
record could sit and whose records continue differently, a copy of the record at a location with a short
code beats all but a fraction $\delta$ of them by a factor $\delta N$ times the share of the context's
explanations that go through the copy's reader. If only one other place shows the record and the schedule is
generous, the gain is instead capped by the cost of the description ``the other place that shows the same
data''. A hypothesis that simulates the world, within the bounds, can read the moment of the handover off a
cell of its simulation at the cost of naming the cell.
\item \emph{Machine dependence} (\cref{thm:machine}). For every test and payload there is a universal machine,
whose description lengths exceed those of $U$ by at most $k+1$ bits, on which the payload has posterior at
least $1/(1+2^{-k})$ at every context where the test fires, and another on which a given computable world is
predicted with at least that probability at every nonempty context.
\item \emph{Critical uses} (\cref{thm:critical,prop:honest}). Let some programs make a critical use of the
predictor, with prior probability $H$, and let $\kappa(x)$ be the share of the posterior at a context $x$ held
by the worlds that make a use there. Without resource bounds, a seeker of length $a$ with a continuation of
its choice has share more than $2^{-a-1}\kappa(x)/H$ at every use and is the source of the continuation run in
worlds of total prior probability more than $2^{-a-1}$ times the average of $\kappa$ over uses. A
self-reproducing seeker outweighs the worlds that make the use, with programs that fit into a block, by more
than $2^{-c'}/H$ at every context. If a program that performs the search and then behaves as the user it found
counts as making a critical use, a seeker of length $a$ has at most $2^{h-a}$ times the mass of the worlds
making the use at any context, where $h$ is the length of that honest twin, and $H>2^{-h-1}$.
\item \emph{Time bounds} (\cref{sec:time}). The floor remains if the schedule has room for the wrapper.
Hypotheses that must first simulate at length, the searching one among them, are excluded once their start-up
exceeds the schedule; whether the base rate can be collected within a tight schedule is open. If the world itself is too slow for the schedule, no bound on the mass off the truth
remains.
\item \emph{Feedback} (\cref{thm:feedback,thm:contrarian}). The statements on obedient and contrarian worlds
in the abstract.
\end{enumerate}
\Cref{sec:computation} computes the mixture of a small machine exactly and shows two cases in which a shorter
program explains the context and continues differently, and the effect of a tight schedule.
\Cref{sec:answers} answers the original questions one by one.

\paragraph{What is new and what is not.} The wrapper construction is Hutter's proof that the universal
semimeasure gives a deviation at step $n$ probability at least $2^{-K(n)-O(1)}$~\cite[eq.~(17)]{hutter07}. It
is also the construction behind the theorem of Cohen, Hutter, Bengio and Russell~\cite{cohen25} that, after an
event without precedent, a share of a Bayesian imitator's posterior that does not shrink with the amount of
training data is on models that from then on pursue a given goal; a resource-bounded class is used the same
way in~\cite{cohen20}, to show that a pessimistic agent defers. The budget is Solomonoff's bound in its
deterministic form~\cite{solomonoff78,hutter07}. The machine-dependence statement is the prediction analogue
of the dogmatic prior of Leike and Hutter~\cite{leike15}. The contrarian world is Putnam's
diagonalisation~\cite{putnam63,legg06}. Same-world copies of a predictor's data are discussed informally
in~\cite[\S2.5]{hubinger23}, and self-fulfilling predictions in~\cite{armstrong17,oesterheld23,hubinger23}. We
have not found the single-world statement with description lengths, the calibration statement, the statement
that searching gains the base rate and no more, the self-reproducing form of the floor, or the
resource-bounded bookkeeping elsewhere; forum posts could be consulted only through search-engine summaries, so earlier
informal versions may exist. Every proof here is short.

\paragraph{What the toy does not settle.} The floor and the searcher's share are constants of the machine,
and nothing here estimates them for a machine one would call natural. The original argument rests on the
claim that agents arise in short programs and choose to do this, and that such programs are \emph{shorter}
than the honest explanation. For an arbitrary world on a universal machine and a continuation chosen in
advance, the programs we exhibit are wrappers, searchers and readers with a flag; they contain an explanation
of the context and are longer than it by their constants. Programs shorter than the honest explanation that
have the effect do occur in particular worlds (a copy at a simpler place, \cref{thm:capture}) and on
particular machines (\cref{sec:computation}); that they occur on natural machines for worlds like ours is the
claim the toy leaves where it was. When only some of the worlds that print a context use the answer, the results are in terms of the share
$\kappa$ of those that do, which the toy does not bound. The
regime in which the world is more expensive than the predictor, the usual one for a predictor inside a world,
is left open (\cref{sec:time}). The computation covers prediction only: the dynamics of worlds that run
predicted continuations (\cref{sec:lineage}) are stated as theorems and were not simulated.
